% Paquetes
\IfPackageLoadedTF{booktabs}{}{\usepackage{booktabs}}
\IfPackageLoadedTF{amsthm}{}{\usepackage{amsthm}}
\IfPackageLoadedTF{longtable}{}{\usepackage{longtable}}
\IfPackageLoadedTF{subfig}{}{\usepackage{subfig}}
% Carga `setspace' para separación entre lineas (si no fue cargado)
\IfPackageLoadedTF{setspace}{}{\usepackage{setspace}}
% Carga `graphicx' si no fue cargado (dado que no se usa entorno pandoc para figures)
\IfPackageLoadedTF{graphicx}{}{\usepackage{graphicx}}
% Carga `etoolbox' 
\IfPackageLoadedTF{etoolbox}{}{\usepackage{etoolbox}}
% mejora la interfaz para objetos flotantes
\IfPackageLoadedTF{float}{}{\usepackage{float}}
% \usepackage{lastpage] % referencia la páginas en su documento LaTex a través de la introducción de una nueva etiqueta

% Definir estilo de citas largas (achica el tamaño)
\IfPackageLoadedTF{relsize}{}{\usepackage{relsize}}
\AtBeginEnvironment{quote}{\smaller}% Step font down one size relative to current font.

% Definir título de tablas/figuras
% \@ifundefined{KOMAClassName}{% if non-KOMA class
% \renewcommand\spanishtablename{Tabla}
%  \addto\captionsspanish{\renewcommand{\tablename}{Tabla}}
%}{% if KOMA class
\renewcaptionname{spanish}{\tablename}{Tabla}
\renewcaptionname{spanish}{\listtablename}{Índice de Tablas}
%}

% Agregado de \source y \notes para titulos secundarios de las tablas y figuras
\usepackage{caption}
\captionsetup{format=plain,labelsep=none,singlelinecheck=false,justification=raggedright,
              font={small,singlespacing,sf},labelfont={sc,bf},width=\textwidth}
\newcommand{\source}[1]{\caption*{\footnotesize{\textit{Fuente:} {#1}}}}
\newcommand{\notes}[1]{\caption*{\footnotesize{\textit{Nota:} {#1}}}}

% Notas sucesivas entre capítulos
\counterwithout{footnote}{chapter}

% Margenes de pagina e imagenes
%\setlength{\headheight}{12.7pt}
\renewcommand{\topfraction}{.85}
\renewcommand{\bottomfraction}{.7}
\renewcommand{\textfraction}{.15}
\renewcommand{\floatpagefraction}{.66}
\setcounter{topnumber}{3}
\setcounter{bottomnumber}{3}
\setcounter{totalnumber}{4}

% Numeración romana para pag preliminares
% \pagenumbering{roman}
\pagenumbering{gobble}

% Quiebre de lineas en chunk
\lstset{
  breaklines=true
}

\makeatletter
\def\thm@space@setup{%
  \thm@preskip = 8pt plus 2pt minus 4pt
  \thm@postskip = \thm@preskip
}

% Modificar títulos de capítulos y partes
\renewcommand*{\partformat}{%
  \partname~\thepart\par
}
\renewcommand*{\chapterformat}{%
  \chapapp~\thechapter\par
}

% Define forma de portada
\usepackage{template/latex/TitleTESIS}

% Definir encabezado de páginas
\usepackage{fancyhdr}
\pagestyle{plain}



%%%% ENCABEZADO DE PÁGINA
% Si utiliza clases de Koma-script debe utilizar el paquete 'scrlayer-scrpage'
% Si utiliza las clases estándar debe usar el paquete 'fancyhdr'
\fancyhead[LE, RE, LO, RO]{}
% \fancyhead[LE]{Pablo Serrati - Tesis de Maestría}
% \fancyhead[RO]{\leftmark}
